\chapter{Ponteiros}

Ao longo da apostila, deixamos algumas perguntas em aberto de propósito. No Capítulo 2, perguntamos por que o \texttt{scanf} exige o operador \texttt{\&} antes do nome da variável. No Capítulo 4, vimos uma função \texttt{trocar} que, por mais correta que parecesse, simplesmente não conseguia trocar os valores de duas variáveis de \texttt{main}. No Capítulo 5, descobrimos que o nome de um vetor se comporta como o endereço do seu primeiro elemento, e que, por isso, funções conseguem alterar vetores recebidos como parâmetro.

Todas essas perguntas têm a mesma resposta: \textbf{ponteiros}. Neste capítulo, vamos finalmente entender o que está por trás de cada uma dessas situações — e, ao fazê-lo, dar o primeiro passo da Parte III, dedicada ao gerenciamento de memória.

\section{O que é um Ponteiro?}

Vamos começar pelo problema, e não pela sintaxe. Desde o Capítulo 2, sabemos que toda variável ocupa uma posição na memória do computador. Cada uma dessas posições possui um \textbf{endereço}: um número que identifica, de forma única, onde aquele dado está armazenado.

Até agora, sempre trabalhamos com o \textbf{valor} das variáveis: lemos o valor, alteramos o valor, passamos o valor para funções. Mas existem situações em que precisamos trabalhar com o \textbf{endereço} de um dado, e não apenas com seu conteúdo — por exemplo, quando queremos que uma função altere uma variável que pertence a outra função.

Um \textbf{ponteiro} é, simplesmente, uma variável capaz de armazenar um endereço de memória.

\begin{lstlisting}
int x = 10;
int *p = &x;
\end{lstlisting}

Nesse trecho, \texttt{x} é uma variável comum, que guarda o valor \texttt{10}. Já \texttt{p} é um ponteiro, que guarda o \textbf{endereço} de \texttt{x}. Podemos visualizar essa relação assim:

\begin{center}
\texttt{
\begin{tabular}{cc}
\begin{tabular}{|c|}
\hline
10 \\
\hline
\end{tabular}
&
\begin{tabular}{|c|}
\hline
0x100 \\
\hline
\end{tabular}
\\[2pt]
\textit{x} (endereço 0x100) & \textit{p} (guarda 0x100)
\end{tabular}
}
\end{center}

Os endereços exibidos aqui (como \texttt{0x100}) são ilustrativos: em um programa real, o valor exato depende do sistema operacional e da execução, e muda de uma vez para outra.

\section{Endereço e Operador \texttt{\&}}

Antes de falarmos do ponteiro em si, vale a pena separar o conceito de \textbf{endereço}. Em C, o operador \texttt{\&}, aplicado a uma variável, devolve o endereço onde ela está armazenada:

\begin{lstlisting}
int x = 10;

printf("%d\n", x);            // exibe o valor de x: 10
printf("%p\n", (void *)&x);   // exibe o endereco de x
\end{lstlisting}

O especificador \texttt{\%p} é usado para exibir endereços de memória (geralmente em hexadecimal), e o \textit{cast} para \texttt{(void *)} é a forma recomendada de passar um endereço para \texttt{printf} com esse especificador.

Repare na diferença entre as duas linhas: \texttt{x} é o \textbf{valor} guardado na variável, e \texttt{\&x} é o \textbf{lugar} onde essa variável vive. Essa distinção — valor versus endereço — é a base de tudo o que vem a seguir.

Isso nos permite, enfim, responder à pergunta que deixamos em aberto no Capítulo 2: por que \texttt{scanf("\%d", \&idade)} precisa do \texttt{\&}? Porque o \texttt{scanf} precisa saber \textbf{onde} escrever o valor lido. Se passássemos apenas \texttt{idade}, ele receberia uma cópia do valor atual da variável — e não teria nenhuma forma de alterar a variável original. Ao passar \texttt{\&idade}, entregamos a ele o endereço onde o valor deve ser gravado.

\section{Declaração de Ponteiros e Operador \texttt{*}}

Para declarar um ponteiro, colocamos um asterisco entre o tipo e o nome da variável:

\begin{lstlisting}
int *p;
\end{lstlisting}

Essa linha deve ser lida como: ``\texttt{p} é um ponteiro para \texttt{int}''. O tipo à esquerda do asterisco indica o \textbf{tipo do dado apontado}. Em seguida, podemos fazer o ponteiro guardar o endereço de uma variável existente:

\begin{lstlisting}
int x = 10;
int *p;

p = &x;
\end{lstlisting}

Agora, o mesmo símbolo \texttt{*} aparece com um segundo significado quando usado em uma \textbf{expressão}, e não em uma declaração:

\begin{lstlisting}
int *p;   // declaracao: p e um ponteiro para int

*p = 20;  // expressao: acessa o valor apontado por p
\end{lstlisting}

\begin{errocomum}
Esquecer que o asterisco \texttt{*} tem dois papéis diferentes em C, dependendo de onde aparece:
\begin{lstlisting}
int *p = &x;  // aqui, * faz parte do TIPO: p e um ponteiro para int
*p = 20;      // aqui, * e o operador de desreferenciacao: acessa o valor apontado
\end{lstlisting}
Na declaração, o \texttt{*} participa da definição do \textbf{tipo} do ponteiro. Em uma expressão, ele realiza a \textbf{desreferenciação}: acessa o valor que está no endereço guardado. Misturar esses dois papéis é uma das maiores fontes de confusão para quem está aprendendo ponteiros — vale reler uma linha com \texttt{*} perguntando-se sempre: ``isto é uma declaração, ou uma expressão?''
\end{errocomum}

\section{Ponteiros e Acesso Indireto}

Com um ponteiro guardando o endereço de uma variável, podemos agora acessar essa variável \textbf{indiretamente}, sem nunca mencionar seu nome:

\begin{lstlisting}
int x = 10;
int *p = &x;

printf("%d\n", x);   // 10
printf("%d\n", *p);  // 10: acessando x atraves de p
\end{lstlisting}

Mais importante: podemos também \textbf{modificar} a variável original através do ponteiro:

\begin{lstlisting}
*p = 50;

printf("%d\n", x);   // agora x vale 50
\end{lstlisting}

Visualmente, o que acontece é:

\begin{center}
\texttt{
\begin{tabular}{ll}
\textbf{Antes:} & x = 10 \quad p $\longrightarrow$ x \\
\textbf{Depois de *p = 50:} & x = 50 \quad p $\longrightarrow$ x \\
\end{tabular}
}
\end{center}

O ponteiro \texttt{p} continua apontando para o mesmo lugar; o que mudou foi o \textbf{conteúdo} desse lugar. A ideia central desta seção pode ser resumida assim:

\begin{quote}
\textit{O ponteiro permite acessar uma variável indiretamente, por meio de seu endereço.}
\end{quote}

\section{Ponteiros e Funções}

Chegamos a uma das seções mais importantes do capítulo, porque ela resolve, finalmente, o problema que deixamos aberto no Capítulo 4. Lembre-se da limitação da passagem por valor:

\begin{lstlisting}
void alterar(int x) {
    x = 100;
}

int main(void) {
    int x = 10;
    alterar(x);
    printf("%d\n", x); // continua 10
    return 0;
}
\end{lstlisting}

A função \texttt{alterar} recebe uma \textbf{cópia} de \texttt{x}, e é essa cópia que é modificada. A variável original permanece intacta.

Para que uma função consiga alterar uma variável de fora, ela precisa receber \textbf{o endereço} dessa variável, e não o seu valor:

\begin{lstlisting}
void alterar(int *x) {
    *x = 100;
}

int main(void) {
    int x = 10;
    alterar(&x);
    printf("%d\n", x); // agora x vale 100
    return 0;
}
\end{lstlisting}

Agora o parâmetro \texttt{x} da função é um ponteiro: ele recebe uma cópia do \textbf{endereço} de \texttt{x}, e, usando \texttt{*x}, a função consegue acessar e modificar o objeto original, que vive em \texttt{main}.

Isso nos permite, enfim, escrever a função \texttt{trocar} do Capítulo 4 de forma correta:

\begin{lstlisting}
void trocar(int *a, int *b) {
    int temp = *a;
    *a = *b;
    *b = temp;
}

// chamada: trocar(&x, &y);
\end{lstlisting}

Há uma sutileza conceitual importante aqui. Em linguagens como C++, existe um mecanismo chamado ``passagem por referência''. \textbf{C não possui passagem por referência propriamente dita}: a passagem continua sendo, sempre, por valor. O que acontece é que o valor passado é um \textbf{endereço} — a função recebe uma cópia do endereço e, por meio dela, consegue modificar o objeto original. O mecanismo de passagem de argumentos não mudou; mudou apenas o que escolhemos passar.

Isso também explica, de forma definitiva, o comportamento de vetores em funções que observamos no Capítulo 5: quando passamos um vetor, estamos passando o endereço do seu primeiro elemento. É por isso que alterações feitas dentro da função afetam o vetor original.

\section{Ponteiros e Vetores}

Agora podemos tornar explícita a relação entre ponteiros e vetores, que já apareceu de forma intuitiva no Capítulo 5. Considere:

\begin{lstlisting}
int v[5] = {10, 20, 30, 40, 50};

int *p = v;
\end{lstlisting}

Como o nome de um vetor é tratado como o endereço do seu primeiro elemento, \texttt{p} passa a apontar para \texttt{v[0]}:

\begin{center}
\texttt{
\begin{tabular}{|c|c|c|c|c|}
\hline
10 & 20 & 30 & 40 & 50 \\
\hline
\end{tabular}
}
\\[2pt]
\texttt{p} aponta para o primeiro elemento (\texttt{v[0]})
\end{center}

A partir daí, podemos acessar os elementos do vetor através do ponteiro:

\begin{lstlisting}
*p        // 10 (primeiro elemento)
*(p + 1)  // 20 (segundo elemento)
*(p + 2)  // 30 (terceiro elemento)
\end{lstlisting}

Aqui aparece, pela primeira vez, a ideia de \textbf{aritmética de ponteiros}: somar um número a um ponteiro faz com que ele ``ande'' pela memória. Vamos formalizar essa ideia na próxima seção.

\section{Aritmética de Ponteiros}

Ponteiros aceitam um conjunto restrito de operações aritméticas:

\begin{lstlisting}
p + 1
p - 1
p++
p--
\end{lstlisting}

O ponto essencial é entender o que significa, de fato, somar um número a um ponteiro. Considere um ponteiro \texttt{int *p}. A expressão \texttt{p + 1} \textbf{não} significa ``o endereço atual mais 1 byte''. Significa avançar \textbf{uma posição do tipo apontado}. Se \texttt{sizeof(int)} for 4, temos:

\begin{center}
\texttt{
\begin{tabular}{|l|l|}
\hline
\textbf{Expressão} & \textbf{Endereço (exemplo)} \\
\hline
p & 0x1000 \\
p + 1 & 0x1004 \\
p + 2 & 0x1008 \\
\hline
\end{tabular}
}
\end{center}

O compilador cuida de multiplicar o deslocamento pelo tamanho do tipo apontado. Por isso, um ponteiro para \texttt{double} avançaria 8 bytes a cada \texttt{+ 1}, e um ponteiro para \texttt{char}, apenas 1.

Essa observação também esclarece a relação entre a notação de vetores e ponteiros. As duas formas abaixo são equivalentes:

\begin{lstlisting}
v[i]
*(v + i)
\end{lstlisting}

Ou seja: o operador de índice \texttt{[]} é, na prática, uma forma mais legível de escrever aritmética de ponteiros seguida de desreferenciação. Esse é um ótimo momento para perceber o quanto a sintaxe de vetores em C está profundamente ligada a ponteiros.

Um cuidado importante: a aritmética de ponteiros só tem comportamento \textbf{bem definido} quando o ponteiro permanece dentro de um mesmo vetor (ou, no máximo, uma posição além do seu final). Somar ou comparar ponteiros que apontam para variáveis não relacionadas — por exemplo, achando que duas variáveis declaradas lado a lado necessariamente ocupam posições vizinhas na memória — não tem garantia nenhuma no padrão da linguagem.

\section{Ponteiros para Diferentes Tipos}

Um ponteiro sempre carrega, junto com o endereço, a informação sobre \textbf{o tipo do dado apontado}. Podemos declarar ponteiros para qualquer tipo que já conhecemos:

\begin{lstlisting}
char   *pc;
int    *pi;
double *pd;
\end{lstlisting}

O tipo apontado importa por três motivos:

\begin{itemize}
    \item determina como os bytes no endereço apontado devem ser \textbf{interpretados} (os mesmos bytes podem significar coisas diferentes como \texttt{int} ou como \texttt{float});
    \item determina o tamanho do \textbf{passo} na aritmética de ponteiros, como vimos na seção anterior;
    \item determina \textbf{quantos bytes} são lidos ou escritos ao desreferenciar o ponteiro (\texttt{*pc} acessa 1 byte, enquanto \texttt{*pd} acessa 8, em sistemas típicos).
\end{itemize}

Isso nos devolve à ideia que apresentamos no Capítulo 2: \textit{o tipo determina como os bits devem ser interpretados}. Um ponteiro apenas estende essa ideia para endereços.

\section{Ponteiros Nulos}

Nem todo ponteiro, em todo momento, aponta para algo válido. Para sinalizar explicitamente que um ponteiro \textbf{não aponta para nenhum objeto}, C oferece a constante \texttt{NULL}:

\begin{lstlisting}
int *p = NULL;
\end{lstlisting}

Conceitualmente, \texttt{NULL} representa a ausência de um endereço válido. Embora \texttt{NULL} costume ser implementado, na prática, como o valor zero, vale pensar nele como um \textbf{conceito} distinto de ``o número zero'': ele diz que o ponteiro, intencionalmente, não aponta para lugar nenhum.

Desreferenciar um ponteiro nulo é um erro, geralmente resultando no encerramento abrupto do programa. Por isso, é comum verificar o ponteiro antes de usá-lo:

\begin{lstlisting}
if (p != NULL) {
    printf("%d\n", *p);
}
\end{lstlisting}

Essa disciplina — inicializar ponteiros com \texttt{NULL} quando ainda não há um endereço válido, e verificar antes de desreferenciar — vai se tornar essencial mais adiante (no Capítulo 10), quando começarmos a solicitar memória dinamicamente durante a execução do programa.

\section{Ponteiros e \texttt{const}}

A palavra-chave \texttt{const}, que apresentamos no Capítulo 2, também pode ser combinada com ponteiros, e o resultado depende de \textbf{onde} ela é escrita. Existem dois efeitos possíveis: proteger o \textbf{valor apontado}, ou proteger o \textbf{próprio ponteiro}.

\begin{lstlisting}
int x = 10;

const int *p = &x;   // ponteiro para const int
\end{lstlisting}

Aqui, o \texttt{const} protege o \textbf{valor apontado}: não é permitido modificar \texttt{x} através de \texttt{p} (\texttt{*p = 20;} seria um erro de compilação), mas o próprio ponteiro \texttt{p} pode ser redirecionado para apontar para outra variável.

\begin{lstlisting}
int *const q = &x;   // ponteiro constante para int
\end{lstlisting}

Aqui, o \texttt{const} protege o \textbf{próprio ponteiro}: \texttt{q} não pode ser redirecionado para outro endereço, mas o valor apontado pode ser modificado normalmente (\texttt{*q = 20;} é válido).

Existe ainda a combinação dos dois efeitos, \texttt{const int *const r = \&x;}, em que nem o ponteiro nem o valor apontado podem ser alterados. Não vamos explorar mais variações aqui: o objetivo é apenas reconhecer que o \texttt{const} pode proteger \textbf{o dado}, \textbf{o ponteiro}, ou \textbf{ambos}. Esse conhecimento já nos dá, inclusive, uma explicação para algo que vimos no Capítulo 6: por que declarar parâmetros de string como \texttt{const char *} sinaliza que a função não vai modificar o texto recebido.

\section{Erros Comuns com Ponteiros}

Ponteiros são uma das fontes mais frequentes de erros em programas C, em parte porque o compilador, na maioria das vezes, não os impede. Estes são os mais comuns:

\begin{errocomum}
\textbf{Ponteiro não inicializado.} Um ponteiro declarado e nunca atribuído contém ``lixo de memória'', exatamente como qualquer variável local:
\begin{lstlisting}
int *p;
*p = 10;  // erro: p aponta para um lugar desconhecido
\end{lstlisting}
Escrever através de um ponteiro cujo endereço é desconhecido pode corromper outras partes da memória do programa, ou encerrá-lo abruptamente.
\end{errocomum}

\begin{errocomum}
\textbf{Desreferenciar \texttt{NULL}.}
\begin{lstlisting}
int *p = NULL;
printf("%d", *p);  // erro: p nao aponta para nada valido
\end{lstlisting}
Sempre verifique se o ponteiro é válido antes de desreferenciá-lo, especialmente quando ele pode ter sido recebido de outra função.
\end{errocomum}

\begin{errocomum}
\textbf{Confundir endereço e valor.}
\begin{lstlisting}
int x = 10;
int *p = x;     // errado: atribui um int a um ponteiro
int *q = &x;    // correto: atribui o endereco de x
\end{lstlisting}
O compilador costuma emitir um aviso nesse caso, mas nem sempre trata como erro. Leve os \textit{warnings} a sério — para isso, vale compilar com as opções \texttt{-Wall -Wextra} no GCC, que ativam avisos extras sobre construções suspeitas.
\end{errocomum}

\begin{errocomum}
\textbf{Confundir \texttt{p} e \texttt{*p}.}
\begin{lstlisting}
p    // o ENDERECO guardado no ponteiro
*p   // o VALOR que esta naquele endereco
\end{lstlisting}
Essa é a confusão mais fundamental, da qual derivam muitas outras: lembre-se sempre de que \texttt{p} e \texttt{*p} são coisas diferentes, de tipos diferentes.
\end{errocomum}

\section*{Resumo do Capítulo}

Podemos condensar tudo o que vimos neste capítulo em uma única cadeia de ideias:

\begin{center}
\textbf{variável $\rightarrow$ valor $\rightarrow$ endereço $\rightarrow$ ponteiro} \\
\textbf{$\rightarrow$ desreferenciação $\rightarrow$ acesso indireto.}
\end{center}

Uma variável ocupa um lugar na memória e guarda um valor. Esse lugar tem um endereço. Um ponteiro é uma variável que guarda esse endereço, e a desreferenciação nos permite acessar, a partir dele, o valor original — de forma indireta. Com essa ideia, resolvemos as perguntas que deixamos em aberto nos capítulos anteriores: o \texttt{\&} do \texttt{scanf}, a função \texttt{trocar} que não funcionava e o comportamento diferente de vetores em funções.

Vimos também como declarar ponteiros, aplicar aritmética de ponteiros dentro de vetores, inicializar ponteiros sem destino com \texttt{NULL} e proteger dados ou endereços com \texttt{const}. Até aqui, todos os ponteiros que utilizamos apontam para objetos cuja memória já existia (variáveis locais ou vetores de tamanho fixo). No próximo capítulo, veremos como a memória de um programa em C é organizada em regiões distintas — preparando o terreno para, no Capítulo 10, aprendermos a solicitar e liberar memória dinamicamente durante a execução.

\begin{objetivos}
\begin{itemize}
    \item Explicar o que é um ponteiro e como ele se relaciona com endereços de memória.
    \item Utilizar os operadores \texttt{\&} e \texttt{*} corretamente, distinguindo seus papéis na declaração e na expressão.
    \item Acessar e modificar uma variável indiretamente por meio de um ponteiro.
    \item Escrever funções que alteram variáveis de quem as chamou, usando ponteiros como parâmetros.
    \item Explicar a relação entre vetores e ponteiros, e utilizar aritmética de ponteiros dentro de um vetor.
    \item Utilizar \texttt{NULL} e verificar ponteiros antes de desreferenciá-los.
    \item Diferenciar \texttt{const int *p} de \texttt{int *const p}.
    \item Reconhecer e evitar os erros mais comuns ao trabalhar com ponteiros.
\end{itemize}
\end{objetivos}

\section*{Exercícios}

\begin{exercicio}

\subsection*{Exercício 1 --- Fixação}
Declare uma variável \texttt{int} com um valor de sua escolha e um ponteiro para ela. Utilize \texttt{printf} para exibir o valor da variável, o endereço da variável (com \texttt{\%p}), o valor do ponteiro e o valor apontado pelo ponteiro. O que você observa ao comparar o endereço da variável com o valor do ponteiro?

\subsection*{Exercício 2 --- Aplicação}
Declare uma variável \texttt{double} e um ponteiro para ela. Modifique o valor da variável \textbf{exclusivamente} através do ponteiro, e confirme, com \texttt{printf}, que a variável original foi alterada.

\subsection*{Exercício 3 --- Funções}
Implemente a função \texttt{void trocar(int *a, int *b)} e utilize-a em um programa de teste. Em seguida, volte ao Exercício 3 do Capítulo 4 e explique, com suas palavras, o que precisou mudar para que a troca passasse a funcionar.

\subsection*{Exercício 4 --- Vetores}
Escreva um programa que declare um vetor de 8 inteiros e o percorra \textbf{utilizando apenas um ponteiro}, sem usar o operador de índice \texttt{[]}, exibindo cada elemento e o endereço em que ele está armazenado. Observe a diferença entre endereços consecutivos e relacione-a com \texttt{sizeof(int)}.

\subsection*{Exercício 5 --- Integrador}
Escreva uma função \texttt{void estatisticas(const int *v, int n, double *media, int *maior, int *menor)} que receba um vetor e seu tamanho e devolva, por meio dos três últimos parâmetros, a média, o maior e o menor elemento. Dentro da função, percorra o vetor utilizando aritmética de ponteiros. Em \texttt{main}, chame a função e exiba os resultados.

\subsection*{Exercício 6 --- Desafio}
O programa abaixo contém \textbf{cinco erros} relacionados a ponteiros, de naturezas diferentes. Encontre-os, explique por que cada um é um problema e corrija o programa.

\begin{lstlisting}
#include <stdio.h>

int main(void) {
    int a = 10;
    int b = 20;
    int *p;
    int *q = a;
    int *r = NULL;

    *p = 5;
    printf("%d\n", *r);
    printf("%d\n", *(&a + 1));

    q = &b;
    printf("%d\n", q);
    return 0;
}
\end{lstlisting}

Depois de corrigi-lo, escreva uma versão em que o último \texttt{printf} exiba, corretamente, o valor de \texttt{b} em vez de seu endereço.

\end{exercicio}